\chapter{Introduction}
\label{chapter:Introduction}

The


\section{Physics Motivation}
\label{sec:intro_physics_motivation}

\subsection{Multi-top quark production at the LHC}
Four-top-quark production (tttt) is one of the rarest Standard Model processes currently accessible at the LHC, with a small but non-zero predicted cross-section.
Its observation provides a direct handle on the top-quark Yukawa coupling, and deviations from the SM prediction would signal new physics beyond the Standard Model.

\subsection{Why classification matters}
Extracting a signal like tttt from its backgrounds (ttH, ttW, ttZ, ttWW) requires powerful event-level discrimination.
Even modest improvements in classifier performance translate into meaningful gains in signal significance at fixed luminosity.

\cite{collaborationObservationFourtopquarkProduction2023}

\section{Machine Learning Challenge}
\label{sec:intro_ml_challenge}

\subsection{Particle events as variable-length sets}
A collision event is a collection of reconstructed objects --- jets, leptons, and missing transverse energy --- with no natural ordering and varying multiplicity per event.
This set-structured data is awkward for architectures that assume fixed-size inputs or sequential ordering.

\subsection{Transformers as a natural fit}
Self-attention is permutation-equivariant by construction, making transformers well-suited to set-valued inputs.
The attention mechanism can in principle learn to weight particle pairs by their physics relevance without any explicit ordering assumption.

\section{Thesis Statement}
\label{sec:intro_thesis_statement}

\subsection{From single model to configurable workbench}
Rather than producing one optimised classifier, this thesis builds a single modular codebase in which every meaningful architectural decision is exposed as a switchable axis.
We introduce 33 orthogonal architectural axes as the organising principle of the entire experimental programme.

\subsection{Controlled ablation as scientific method}
By varying one axis at a time while holding all others fixed, each experiment becomes a controlled ablation.
The central question is not which configuration achieves the highest AUROC, but what architectural choices matter, and why.

\section{Deliverables}
\label{sec:intro_deliverables}

\subsection{Modular codebase}
The primary engineering output is a Hydra-driven Python package in which every architecture variant is reproducible from a single configuration file.
Any run appearing in this thesis can be regenerated with a single CLI command.

\subsection{Experiment database}
A Weights \& Biases project serves as a structured record of all training runs, storing hyperparameters, metrics, and axis settings in a queryable format.
This database is the raw evidence underlying every quantitative claim in the thesis.

\subsection{Thesis document}
This document is a curated, interpretable distillation of the experiment database: not an exhaustive log of all runs, but a set of focused scientific arguments each anchored to a specific axis group.

\section{Reader's Guide}
\label{sec:intro_readers_guide}

\subsection{Structure of the thesis}
The thesis is organised in three parts: Part~I frames the problem and the workbench (Chapters~2--3); Part~II addresses one focused scientific question per chapter, each mapped to a specific axis group (Chapters~4--9); Part~III synthesises the findings into a recommended configuration and outlook (Chapters~10--11).

\subsection{How to read the results}
Throughout the thesis, AUROC is the primary metric, reported as a seed-averaged mean with error bars.
Axis groups are referred to by their labels (D, A, B, P, H), which correspond directly to config namespaces in the codebase.
Each Part~II chapter is framed by one of three evaluation lenses --- performance, interpretability, or efficiency --- which is stated explicitly at the chapter opening.
